% ── 2. Related Work ───────────────────────────────────────────────────────────
\section{Related Work}
\label{sec:related}

\textbf{LightRAG}~\cite{guo2024lightrag} introduces dual-level retrieval:
a low-level mode that expands the 1-hop neighborhood of entity anchors, and
a high-level mode that traverses relation embeddings.
Mode selection uses coarse keyword signals, not query-shape classification,
and both modes apply the same fixed-radius expansion regardless of whether
the question asks for a definition, a bridge traversal, or a thematic sweep.

\textbf{HippoRAG}~\cite{gutierrez2024hipporag} seeds Personalized PageRank
(PPR)~\cite{page1999pagerank} from query entities, letting relevance diffuse
through the graph in a way that naturally captures multi-hop associations.
The limitation is uniformity: a single-entity definitional question receives
the same global diffusion as a question asking for synthesis across multiple
anchors.
PPR is well-suited to thematic queries; it is a poor fit for questions that
require a precise path between two named entities.

\textbf{G-Retriever}~\cite{he2024g} computes a prize-collecting Steiner tree
spanning all query entities, which is effective for multi-entity questions
but degenerates to a single node for single-entity queries---the opposite
failure mode from HippoRAG.

\textbf{PathRAG}~\cite{chen2025pathrag} is the closest prior work.
It extracts relational paths between anchor pairs using flow-based pruning
and presents them in reliability order, exploiting LLM recency bias to
surface the most relevant path at the end of the context.
This works well on bridge and relational-chain questions.
The gap is that PathRAG's path format is the same for every query type: a
definitional question gets a list of paths from the anchor to nearby
entities, a thematic question gets paths between whichever entities the
entity extractor picks.
The key thing GRASP-RAG adds is \emph{selecting the retrieval motif and
context format conditional on the query's logical shape}.

\textbf{Layered Query Retrieval (LQR)}~\cite{huang2024layered} shares
GRASP-RAG's core motivation---query properties should govern retrieval
strategy---and classifies each query by complexity (binary: simple vs.\
complex).
Simple queries get a single BM25 retrieval pass; complex queries are
decomposed into two sub-questions for iterative retrieval, followed by an
NLI relevance filter.
The binary split, however, cannot distinguish between an ego-network query
and a bridge traversal query---both would be labelled ``complex'' and receive
the same decomposition treatment---and LQR operates on flat text without a
knowledge graph.

\textbf{BridgeRAG}~\cite{bacellar2026bridgerag} addresses the same failure
mode as our bridge motif from a graph-free angle.
It conditions second-hop candidate scoring on an intermediate bridge passage
via a tripartite judge $s(q, b, c)$, achieving strong recall on HotpotQA,
2WikiMultiHopQA, and MuSiQue.
Unlike GRASP-RAG, BridgeRAG has no knowledge graph and applies bridge
conditioning uniformly to all queries; GRASP-RAG reserves bridge-path
traversal for queries whose linguistic form signals a relational chain, while
routing definitional and thematic questions to structurally matched motifs.
